% 
Remark: Numbers in tiles indicate the heuristic of the neighbor (if the queen in this file is moved to this tile).
\begin{enumerate}[(a)]
    \item Hill Climbing with heuristic and neighborhood from slide C1 p. 23.\\
    Initial candidate (current heuristic value: 3):
\begin{center}
  \begin{tikzpicture}
    \draw[step=1,black,thin] (0,0) grid (5,5);

    % First file
    \node[BlackSquare] at (0,0) {3}; 
    \node[WhiteSquare] at (0,1) {3};
    \node[BlackSquare] at (0,2) {\BlackQueenOnWhite};
    \node[WhiteSquare] at (0,3) {2};
    \node[BlackSquare] at (0,4) {\color{blue}2\color{black}};
    
    % Second file   
    \node[WhiteSquare] at (1,0) {4};
    \node[BlackSquare] at (1,1) {\BlackQueenOnWhite};
    \node[WhiteSquare] at (1,2) {5};
    \node[BlackSquare] at (1,3) {5};
    \node[WhiteSquare] at (1,4) {4};

    % Third file 
    \node[BlackSquare] at (2,0) {5};
    \node[WhiteSquare] at (2,1) {4};
    \node[BlackSquare] at (2,2) {6};
    \node[WhiteSquare] at (2,3) {\BlackQueenOnWhite};
    \node[BlackSquare] at (2,4) {3};

    % Fourth file
    \node[WhiteSquare] at (3,0) {2};
    \node[BlackSquare] at (3,1) {3};
    \node[WhiteSquare] at (3,2) {\BlackQueenOnWhite};
    \node[BlackSquare] at (3,3) {3};
    \node[WhiteSquare] at (3,4) {2};

    % Fifth file
    \node[BlackSquare] at (4,0) {\BlackQueenOnWhite};
    \node[WhiteSquare] at (4,1) {6};
    \node[BlackSquare] at (4,2) {5};
    \node[WhiteSquare] at (4,3) {5};
    \node[BlackSquare] at (4,4) {4};

    \end{tikzpicture}
\end{center}
First iteration (current heuristic value: 2):
\begin{center}
  \begin{tikzpicture}
    \draw[step=1,black,thin] (0,0) grid (5,5);

    % First file
    \node[BlackSquare] at (0,0) {3}; 
    \node[WhiteSquare] at (0,1) {3};
    \node[BlackSquare] at (0,2) {3};
    \node[WhiteSquare] at (0,3) {2};
    \node[BlackSquare] at (0,4) {\BlackQueenOnWhite};
    
    % Second file   
    \node[WhiteSquare] at (1,0) {4};
    \node[BlackSquare] at (1,1) {\BlackQueenOnWhite};
    \node[WhiteSquare] at (1,2) {4};
    \node[BlackSquare] at (1,3) {5};
    \node[WhiteSquare] at (1,4) {5};

    % Third file 
    \node[BlackSquare] at (2,0) {3};
    \node[WhiteSquare] at (2,1) {3};
    \node[BlackSquare] at (2,2) {5};
    \node[WhiteSquare] at (2,3) {\BlackQueenOnWhite};
    \node[BlackSquare] at (2,4) {2};

    % Fourth file
    \node[WhiteSquare] at (3,0) {2};
    \node[BlackSquare] at (3,1) {4};
    \node[WhiteSquare] at (3,2) {\BlackQueenOnWhite};
    \node[BlackSquare] at (3,3) {3};
    \node[WhiteSquare] at (3,4) {3};

    % Fifth file
    \node[BlackSquare] at (4,0) {\BlackQueenOnWhite};
    \node[WhiteSquare] at (4,1) {4};
    \node[BlackSquare] at (4,2) {2};
    \node[WhiteSquare] at (4,3) {3};
    \node[BlackSquare] at (4,4) {3};

    \end{tikzpicture}
\end{center}
The algorithm already stops after the first iteration, because the heuristic value
of the current candidate is 2 and there is no tile with a better heuristic value.

\newpage
    \item Hill climbing variant where the next candidate that is considered is selected
by first picking a file at random and then moving the (unique) queen in that file to a square
in that file with the best heuristic value (the variant is introduced on slide 9 of Chapter C2
in the print version of the lecture slides). Assume that the random choices are resolved in
such a way that the file selection picks the third file first, and in each subsequent step the file
that is to the left of the file that was selected last (if the leftmost file was selected last, the
rightmost file is selected).\\[0.5em]
\color{red}TODO \color{black}
\end{enumerate}